\section*{Capítulo \ref{ch:g12:kinematics-2d} --- Cinemática em duas dimensões}

\begin{solution}{exo:g12:kinematics-2d:1}
$\vect v = (3.0, 4.0)$, constante; \omterm{def:g12:kinematics-2d:velocity}{velocidade escalar}
$\sqrt{3.0^2 + 4.0^2} = \qty{5.0}{m/s}$. \omterm{def:g10:relative-motion:trajectory}{Trajetória}: a reta
$y = \tfrac43 x$. \omterm{def:g12:kinematics-2d:velocity}{Vetor velocidade} constante: \omterm{def:g12:kinematics-2d:uniform}{movimento retilíneo
uniforme}.
\end{solution}

\begin{solution}{exo:g12:kinematics-2d:2}
Reta, com espaçamentos iguais: \omterm{def:g12:kinematics-2d:uniform}{retilíneo uniforme};
$v = 0.35/0.10 = \qty{3.5}{m/s}$. Espaçamentos que se abrem
significariam \omterm{def:g12:kinematics-2d:velocity}{velocidade escalar} crescente.
\end{solution}

\begin{solution}{exo:g12:kinematics-2d:3}
$v(t) = 4.0\,t$, $a = \qty{4.0}{m/s^2}$: aceleração constante, movimento
retilíneo uniformemente acelerado. $v(3.0) = \qty{12}{m/s}$.
\end{solution}

\begin{solution}{exo:g12:kinematics-2d:4}
$f = 1/8.0 = \qty{0.125}{Hz}$;
$v = 2\pi \times 4.5/8.0 \approx \qty{3.5}{m/s}$;
$a = v^2/R = 3.53^2/4.5 \approx \qty{2.8}{m/s^2}$, apontando para o eixo
do carrossel.
\end{solution}

\begin{solution}{exo:g12:kinematics-2d:5}
Inclinação: $a = (10.0 - 2.0)/4.0 = \qty{2.0}{m/s^2}$. Área (trapézio):
$d = \tfrac12 (2.0 + 10.0) \times 4.0 = \qty{24}{m}$.
\end{solution}

\begin{solution}{exo:g12:kinematics-2d:6}
$v = \qty{88.9}{m/s}$: $a_n = 88.9^2/6000 \approx \qty{1.3}{m/s^2}
\approx 0.13\,g$. Como $a_n = v^2/R$, um $v$ alto exige um $R$ enorme
para manter aceitável a aceleração lateral (e a \omterm{def:g10:inertia:force}{força} sobre a via e sobre
os passageiros) --- a inclinação da via muda quem a fornece, e não o
tamanho dela.
\end{solution}

\begin{solution}{exo:g12:kinematics-2d:7}
$v(t) = -2.4\,t + 8.0$; $a = \qty{-2.4}{m/s^2}$, constante. Inversão em
$v = 0$: $t = 8.0/2.4 \approx \qty{3.3}{s}$, em
$x \approx \qty{16.3}{m}$. Antes: para a frente, freando; depois: para
trás, ganhando velocidade --- uniformemente acelerado o tempo todo.
\end{solution}

\begin{solution}{exo:g12:kinematics-2d:8}
$v_1 = 0.28/0.40 = \qty{0.70}{m/s}$;
$v_2 = (0.54 - 0.10)/0.40 = \qty{1.10}{m/s}$;
$v_3 = (0.88 - 0.28)/0.40 = \qty{1.50}{m/s}$. A velocidade ganha
\qty{0.40}{m/s} a cada \qty{0.20}{s}: $a = \qty{2.0}{m/s^2}$, constante
--- retilíneo uniformemente acelerado.
\end{solution}

\begin{solution}{exo:g12:kinematics-2d:9}
$T = \qty{2.36e6}{s}$:
$v = 2\pi \times \num{3.84e8}/\num{2.36e6} \approx \qty{1.0e3}{m/s}$;
$a_n = v^2/R \approx \qty{2.7e-3}{m/s^2}$. Razão:
$g/a_n \approx 3600 = 60^2$ --- a gravidade, diluída pelo inverso do
quadrado da distância (\cref{ch:g10:universal-gravitation}), explica
exatamente a curva da Lua.
\end{solution}

\begin{solution}{exo:g12:kinematics-2d:10}
$f = 1200/60 = \qty{20}{Hz}$, $T = \qty{0.050}{s}$;
$v = 2\pi R f = 2\pi \times 0.25 \times 20 \approx \qty{31}{m/s}$;
$a_n = v^2/R \approx \qty{3.9e3}{m/s^2} \approx 400\,g$. A parede do
tambor fornece o $a_n$ da roupa; pelos furos nada empurra a água para
dentro, e por isso ela sai pela tangente.
\end{solution}

\begin{solution}{exo:g12:kinematics-2d:11}
$v_0 = \qty{25}{m/s}$: $t_s = 25/7.5 \approx \qty{3.3}{s}$;
$d = v_0^2/(2a) = 625/15 \approx \qty{42}{m}$. Gráfico: um segmento de
reta descendo de \qty{25}{m/s} até $0$ em $t_s$.
\end{solution}

\begin{solution}{exo:g12:kinematics-2d:12}
$\vect v = (-R\omega\sin\omega t,\, R\omega\cos\omega t)$, de norma
constante $R\omega$; o produto escalar dele com $\vect{OM}$ se anula, de
modo que ele é perpendicular ao raio e, portanto, tangente. Derivando de
novo, $\vect a = -\omega^2\vect{OM}$: em direção a $O$, de norma
$\omega^2 R = (v/R)^2 R = v^2/R$.
\end{solution}

\begin{solution}{exo:g12:kinematics-2d:13}
$a_t = \qty{2.5}{m/s^2}$ (para trás),
$a_n = 20^2/80 = \qty{5.0}{m/s^2}$;
$a = \sqrt{2.5^2 + 5.0^2} \approx \qty{5.6}{m/s^2}$; ângulo com a direção
do movimento:
$\theta = 180^\circ - \arctan(5.0/2.5) \approx 117^\circ$ (ou seja,
$63^\circ$ a partir de exatamente para trás). $5.6 < 7.0$: dentro da
aderência, com pouca folga.
\end{solution}

\begin{solution}{exo:g12:kinematics-2d:14}
Arrancada: $t_1 = 20/1.2 \approx \qty{17}{s}$,
$d_1 = 20^2/(2 \times 1.2) \approx \qty{167}{m}$. Cruzeiro:
\qty{40}{s}, $\qty{800}{m}$. Frenagem:
$t_3 = 20/1.5 \approx \qty{13}{s}$,
$d_3 = 400/3.0 \approx \qty{133}{m}$. Total: $\qty{1.1e3}{m}$ em
\qty{70}{s}; média $1100/70 \approx \qty{16}{m/s}$ (\qty{57}{km/h}). O
$v(t)$ é um trapézio cuja área,
$\tfrac12(70 + 40) \times 20 = \qty{1100}{m}$, confere com o total.
\end{solution}

\begin{solution}{exo:g12:kinematics-2d:15}
$v = 2\pi \times \num{4.22e7}/\num{86164} \approx \qty{3.1e3}{m/s}$;
$a_n = v^2/R \approx \qty{0.22}{m/s^2}$. Quem a fornece é a gravidade da
Terra àquela distância; o \omterm{def:g12:kinematics-2d:ucm}{período} é igual ao \omterm{def:g12:kinematics-2d:ucm}{período} de rotação da
Terra, de modo que \omterm{def:g10:universal-gravitation:satellite}{satélite} e solo giram juntos e ele fica pairando sobre
um ponto do equador (\cref{ch:g12:satellites-kepler}).
\end{solution}

\begin{solution}{pb:g12:kinematics-2d:1}
\textbf{1.} O movimento é relativo (\cref{ch:g10:relative-motion}): toda
posição e toda velocidade precisam nomear o seu \omterm{def:g10:relative-motion:frame}{referencial}. Aqui: o
\omterm{def:g10:relative-motion:frames}{referencial do solo} (a estrada), com $x$ ao longo da estrada a partir da
posição do primeiro quadro.

\textbf{2.} Em cada intervalo: $9.0/0.50 = \qty{18.0}{m/s}$ --- quatro
valores iguais.

\textbf{3.} Estrada reta, deslocamentos iguais em tempos iguais:
\omterm{def:g12:kinematics-2d:uniform}{movimento retilíneo uniforme}.

\textbf{4.} $v_0 = \qty{18.0}{m/s} = \qty{64.8}{km/h}$.

\textbf{5.} $\vect a = \vect 0$.

\textbf{6.} $v(t) = -a\,t + v_B$; $x(t) = -\tfrac12 a t^2 + v_B\,t$.

\textbf{7.} $v(t_s) = 0$ dá $t_s = v_B/a$; então
$d_b = -\tfrac12 a (v_B/a)^2 + v_B^2/a = v_B^2/(2a)$.

\textbf{8.} $v_B = \sqrt{2 a d_b} = \sqrt{2 \times 6.0 \times 27} =
\sqrt{324} = \qty{18.0}{m/s}$ --- exatamente a velocidade filmada: o
motorista não reduziu antes de frear.

\textbf{9.} Reação: $18.0 \times 1.0 = \qty{18}{m}$; total
$18 + 27 = \qty{45}{m}$.

\textbf{10.} $d_b = v_B^2/(2a) \propto v_B^2$: dobrar $v_B$ multiplica
$d_b$ por $4$. A \qty{50}{km/h} ($\qty{13.9}{m/s}$):
$d_b = 13.9^2/12.0 \approx \qty{16}{m}$.

\textbf{11.} A \omterm{def:g12:kinematics-2d:velocity}{velocidade escalar} é constante, mas a direção da
velocidade gira: $\vect v$ não é constante, de modo que o carro está
acelerando.

\textbf{12.} Quarto de volta: $2\pi R/4 = \qty{19.6}{m}$;
$v = 19.6/2.5 \approx \qty{7.9}{m/s}$ (\qty{28}{km/h}).

\textbf{13.} $a_n = 7.85^2/12.5 \approx \qty{4.9}{m/s^2}$, em direção ao
centro da rotatória.

\textbf{14.} $a_{\max} = \mu g = 0.70 \times 9.81 \approx
\qty{6.9}{m/s^2}$;
$v_{\max} = \sqrt{\mu g R} = \sqrt{6.87 \times 12.5} \approx
\qty{9.3}{m/s}$ (\qty{33}{km/h}). $7.9 < 9.3$: dentro da aderência ---
coerente com a faixa limpa, sem derrapagem.

\textbf{15.} $T = 2\pi R/v = 78.5/7.85 \approx \qty{10.0}{s};
f = \qty{0.10}{Hz}$.

\textbf{16.} $64.8 - 50 \approx \qty{15}{km/h}$ acima --- cerca de
$30\%$ acima do limite.

\textbf{17.} A frenagem age ao longo de $41 - 18 = \qty{23}{m}$:
$v^2 = v_B^2 - 2ad = 324 - 2 \times 6.0 \times 23 = 48$, logo
$v \approx \qty{6.9}{m/s} \approx \qty{25}{km/h}$ no impacto.

\textbf{18.} A \qty{13.9}{m/s}:
$13.9 \times 1.0 + 13.9^2/12.0 = 13.9 + 16.1 = \qty{30.0}{m} <
\qty{41}{m}$ --- o carro para \qty{11}{m} antes da van.

\textbf{19.} $d(v) = v\,t_r + v^2/(2a)$: uma vez linearmente (a reação
cobre $v t_r$ antes de os freios agirem) e outra ao quadrado (a parcela
cinética $v^2/(2a)$) --- e é por isso que excessos moderados custam de
modo desproporcional.

\textbf{20.} A câmera de bordo e as marcas de derrapagem concordam:
\qty{65}{km/h} numa zona de \qty{50}{km/h}; no limite legal o carro teria
parado \qty{11}{m} antes da van --- o acidente são os \qty{15}{km/h} a
mais.
\end{solution}
